\section{Grazing-uniform endpoint localization and the third variation}
\label{sec:el-localization}
The estimates in this section are on every actual regular finite itinerary,
not on an auxiliary stochastic process. They hold for all $m\geq1$ and
$R\in[9/20,47/100]$, without a lower cutoff on any positive incidence.
The derivative is with respect to the two endpoint angles of the actual
stationary action. It is not the derivative of the forward collision map
with fixed initial momentum. This distinction is essential: no uniform
bound on that expanding forward derivative is asserted.

Keep $\tau_-=3/50$, $R_+=47/100$, and define
\begin{equation}\label{eq:el-constants}
 \sigma=\frac{47}{59},\qquad
 s_i=\sigma^i+\sigma^{m-i},\qquad
 J_* =\max\left\{1,\frac{R_+}{2\tau_-}(\sigma+\sigma^{-1})\right\}
       =\frac{2845}{354}.
\end{equation}
On a local physical branch let
$\alpha(q)=(\alpha_0(q),\ldots,\alpha_m(q))$, $q=(\alpha_0,\alpha_m)$,
be the stationary lift. For an endpoint direction $\eta\in\R^2$ put
$\xi_i=D_q\alpha_i[\eta]$. Every assertion about continuation below is
restricted to a path in this branch up to its first actual boundary.

\begin{theorem}[Two-sided localization without an inverse-incidence loss]
\label{thm:el-jacobi}
For every such trajectory and every endpoint direction,
\begin{align}
 |c_i\xi_i|&\leq
     \sigma^i c_0|\eta_0|+\sigma^{m-i}c_m|\eta_m|
     \leq s_i\|\eta\|_\infty,\label{eq:el-weighted}\\
 |\xi_i|&\leq J_*s_i\|\eta\|_\infty,
                         \qquad 0\leq i\leq m.\label{eq:el-unweighted}
\end{align}
In particular, endpoint displacement decays exponentially towards the
middle of the true collision chain even when an intermediate incidence
is arbitrarily small. No comparison trajectory, new section, or collision
suppression is used.
\end{theorem}
\begin{proof}
For $m=1$ the endpoint assertions follow immediately from $s_0,s_m\geq1$.
Suppose $m\geq2$. Differentiate the interior stationary equations and
use the full physical second variation \eqref{eq:v109-second-variation}.
With $w_i=\tau_i^{-1}$ and $y_i=(-1)^ic_i\xi_i$ this gives
\begin{equation}\label{eq:el-dirichlet}
 (\kappa_i+w_{i-1}+w_i)y_i-w_{i-1}y_{i-1}-w_iy_{i+1}=0,
 \qquad \kappa_i=\frac2{Rc_i},\quad 1\leq i<m.
\end{equation}
The sign change is important: both off-diagonal coefficients in this
Dirichlet equation are negative. For a real sequence with nonnegative
boundary values and nonnegative left side in \eqref{eq:el-dirichlet},
a negative interior minimum would make that left side strictly negative,
because $\kappa_i>0$. This proves the needed maximum principle directly.

Write the left side as $\mathcal D y$. Since $w_i\leq\tau_-^{-1}$ and
$\kappa_i\geq2/R_+$,
\[
 \frac{\mathcal D(\sigma^i)_i}{\sigma^i}
 =\kappa_i+w_{i-1}(1-\sigma^{-1})+w_i(1-\sigma)
 \geq\frac2{R_+}-\frac{\sigma^{-1}-1}{\tau_-}=0.
\]
The reversed sequence $\sigma^{m-i}$ is a supersolution by the same
calculation with the adjacent flights interchanged. Thus
$b_i=\sigma^ic_0|\eta_0|+\sigma^{m-i}c_m|\eta_m|$ is a nonnegative
supersolution and $b_0\geq|y_0|$, $b_m\geq|y_m|$.
Apply the maximum principle separately to $b-y$ and $b+y$ to get
\eqref{eq:el-weighted}.

Dividing that bound by $c_i$ would lose uniformity and is not the next
step. Instead use the equation itself:
\[
 \kappa_i|y_i|
 \leq w_{i-1}|y_{i-1}|+w_i|y_{i+1}|.
\]
Since $\kappa_ic_i=2/R$, it follows that
\[
 |\xi_i|\leq\frac{R_+}{2\tau_-}(s_{i-1}+s_{i+1})\|\eta\|_\infty
 =\frac{R_+}{2\tau_-}(\sigma+\sigma^{-1})s_i\|\eta\|_\infty.
\]
The endpoint inequalities follow from $\xi_0=\eta_0$, $\xi_m=\eta_m$.
This proves \eqref{eq:el-unweighted} for the entire chain. The calculation
used positive incidences only at the actual regular point, with no
uniform positive lower bound on them.
\end{proof}

\subsection{Physical flights, incidence, and every unvisited disk}
Let $v_i=(x_{i+1}-x_i)/\tau_i$ be the actual unit flight direction.
Put
\begin{equation}\label{eq:el-observable-constants}
 C_c=J_*\left[1+\frac{R_+}{\tau_-}(1+\sigma^{-1})\right],\qquad
 C_\Delta=R_+J_*(1+\sigma).
\end{equation}

\begin{proposition}[Differentiation of the physical observables]
\label{prop:el-observables}
On each actual regular branch,
\begin{align}
 |D_q\tau_i[\eta]|&\leq R_+J_*(s_i+s_{i+1})\|\eta\|_\infty,
                                      &&0\leq i<m,\label{eq:el-flight}\\
 \|D_qv_i[\eta]\|_2&\leq\frac{R_+J_*}{\tau_-}
                (s_i+s_{i+1})\|\eta\|_\infty,
                                      &&0\leq i<m,\label{eq:el-direction}\\
 |D_qc_i[\eta]|&\leq C_cs_i\|\eta\|_\infty,
                                      &&0\leq i\leq m.\label{eq:el-incidence}
\end{align}
The true least clearance $\Delta_i$, with precisely the two incident
disks excluded as in \eqref{eq:v76-guard}, is locally Lipschitz and has
the following directional bound along every continuously differentiable
physical branch path, almost everywhere on that path:
\begin{equation}\label{eq:el-clearance}
 |D_q\Delta_i[\eta]|\leq
 R_+J_*\max(s_i,s_{i+1})\|\eta\|_\infty
 \leq 2C_\Delta\sigma^{d'_i}\|\eta\|_\infty.
\end{equation}
The last assertion includes changes of the nearest unvisited disk and
of its nearest point on the flight. There is no factor equal to the
number of candidate disks.
\end{proposition}
\begin{proof}
$D_qx_i[\eta]=Rn'(\alpha_i)\xi_i$. Differentiate length and unit
direction and use respectively $\|D|X|\|=1$ and
$D(X/|X|)=|X|^{-1}(I-vv^{\mathsf T})$ to prove
\eqref{eq:el-flight}--\eqref{eq:el-direction}. For $i<m$,
$c_i=v_i\cdot n_i$, so $|D c_i|\leq\|Dv_i\|_2+|\xi_i|$.
For the terminal collision use $c_m=-v_{m-1}\cdot n_m$ instead.
The inequalities $s_{i+1}\leq\sigma^{-1}s_i$ and
$s_{i-1}\leq\sigma^{-1}s_i$ give \eqref{eq:el-incidence}.
This argument differentiates the actual geometric cosine, not an
expression with an artificially introduced $1/c_i$.

The Hausdorff distance of segments $[x_i,x_{i+1}]$ and
$[\widetilde x_i,\widetilde x_{i+1}]$ is at most
$\max(\|x_i-\widetilde x_i\|_2,\|x_{i+1}-\widetilde x_{i+1}\|_2)$:
compare points with the same affine parameter. Distance to any fixed
center changes by at most that number. Taking an infimum over exactly
the same set of unvisited centers preserves this inequality. Subtracting
$R$ does not change it. The fixed itinerary fixes the two excluded
centers; it does not discard any other disk. This proves the local
Lipschitz assertion and its derivative bound along a branch path by
integration of the endpoint speeds. A Lipschitz function on a compact
path interval is absolutely continuous; nearest-disk switches do not
create atomic jumps. Finally
\[
 s_i+s_{i+1}=(1+\sigma)(\sigma^i+\sigma^{m-1-i})
             \leq2(1+\sigma)\sigma^{d'_i}
\]
proves the last inequality in \eqref{eq:el-clearance}.
\end{proof}

\subsection{A third-variation bound uniform in collision count}
Define constants depending only on the original geometric interval,
\begin{equation}\label{eq:el-third-constant}
 \begin{split}
 C_3&=R_++\frac{3R_+^2}{\tau_-}+\frac{3R_+^3}{\tau_-^2},\\
 C_F&=\frac{8C_3J_*^3(1+\sigma)^3}{1-\sigma^3}.
 \end{split}
\end{equation}

\begin{theorem}[Uniform physical third variation]
\label{thm:el-third}
For every actual regular endpoint action, all $m\geq1$, all original
radii, and $\eta,\zeta,\upsilon\in\R^2$,
\begin{equation}\label{eq:el-third}
 |D_q^3F[\eta,\zeta,\upsilon]|
 \leq C_F\|\eta\|_\infty\|\zeta\|_\infty\|\upsilon\|_\infty.
\end{equation}
Consequently $\|D_qG[\eta]\|_{\rm op}\leq2C_F\|\eta\|_\infty$.
For any piecewise continuously differentiable path $q(a)$ that remains
in a single actual regular branch,
\[
 \|G(q(b))-G(q(a))\|_{\rm op}
 \leq2C_F\int_a^b\|q'(v)\|_\infty\,dv.
\]
In particular this bound applies on every line segment contained in
that branch. It does not assert the existence of a collision-uniform
ball of physical continuation across unvisited tangencies or section cuts.
\end{theorem}
\begin{proof}
For the Euclidean norm $N(X)=|X|$, at $X\ne0$,
$\|DN\|\leq1$, $\|D^2N\|\leq|X|^{-1}$ and
$\|D^3N\|\leq3|X|^{-2}$. The last inequality follows by writing its
trilinear form as minus $|X|^{-2}$ times the sum of
$(v\cdot a)b^{\mathsf T}(I-vv^{\mathsf T})c$ and the two cyclic terms.
For one flight $X=R n(\alpha_{i+1})-R n(\alpha_i)+O_{i+1}-O_i$,
each derivative of $n$ up to order three has norm one. If
$V_i(\xi)=|\xi_i|+|\xi_{i+1}|$, the third-order chain rule therefore gives
\begin{equation}\label{eq:el-flight-third}
 |D^3\tau_i[\xi,\widetilde\xi,\widehat\xi]|
 \leq C_3 V_i(\xi)V_i(\widetilde\xi)V_i(\widehat\xi).
\end{equation}
The coefficient three in the middle term of $C_3$ counts the three
second--first derivative terms; no such term is omitted.

Let $\mathcal J\eta=D_q\alpha[\eta]$ be the stationary lift.
Differentiated stationarity says that $H\mathcal J\eta$ has zero
interior coordinates, where $H=D^2S$. From
$D^2F[\eta,\zeta]=H[\mathcal J\eta,\mathcal J\zeta]$, differentiation in
$\upsilon$ gives the exact identity
\begin{equation}\label{eq:el-third-schur}
 D^3F[\eta,\zeta,\upsilon]
 =D^3S[\mathcal J\eta,\mathcal J\zeta,\mathcal J\upsilon].
\end{equation}
Indeed, derivatives of $\mathcal J$ have zero endpoint coordinates,
and their contractions with $H\mathcal J\eta$ or
$H\mathcal J\zeta$ vanish. This cancellation occurs after differentiating
the full stationary action; it is not a claim that derivatives of the
interior variables vanish.

By Theorem~\ref{thm:el-jacobi}, each $V_i(\mathcal J\eta)$ is at most
$J_*(s_i+s_{i+1})\|\eta\|_\infty$. Moreover,
\[
 \sum_{i=0}^{m-1}(s_i+s_{i+1})^3
 \leq4(1+\sigma)^3\sum_{i=0}^{m-1}
       (\sigma^{3i}+\sigma^{3(m-1-i)})
 \leq\frac{8(1+\sigma)^3}{1-\sigma^3}.
\]
Sum \eqref{eq:el-flight-third} and use \eqref{eq:el-third-schur} to
obtain \eqref{eq:el-third}. Each entry of $D_qG[\eta]$ is bounded by
$C_F\|\eta\|_\infty$; its operator norm is at most twice that number.
Integrating along the stated physical path completes the proof.
\end{proof}

\begin{remark}[Domain, time, and source]
\label{rem:el-domain}
Theorem~\ref{thm:el-third} concerns the true reduced action at every
regular point. It requires neither a pressure theorem nor the
endpoint-momentum kernel. Uniform boundedness of its third derivative
is not uniform control of all higher derivatives, a chart-count bound,
or a roof-density theorem. All labels may restrict the points at which
it is used, but are not differentiated. For a path of genuine trajectories
with a fixed itinerary reaching a physical boundary the estimates hold
up to that boundary. They do not define a collision-map derivative at a
visited tangency or permit continuing an exact event through a new hit.
\end{remark}
